\documentclass[10pt,a4paper]{amsart}
\usepackage{amsmath,amssymb,amsthm,graphicx}
\usepackage{verbatim}
\newenvironment{reference}{\comment}{\endcomment}
%\newenvironment{reference}{}{}
\newenvironment{slogan}{\comment}{\endcomment}
\newenvironment{history}{\comment}{\endcomment}

% For commutative diagrams we use Xy-pic
\usepackage[all]{xy}

% We use 2cell for 2-commutative diagrams.
\xyoption{2cell}
% Theorem environments.
%
\theoremstyle{plain}
\newtheorem{theorem}[subsection]{Theorem}
\newtheorem{proposition}[subsection]{Proposition}
\newtheorem{lemma}[subsection]{Lemma}

\theoremstyle{definition}
\newtheorem{definition}[subsection]{Definition}
\newtheorem{example}[subsection]{Example}
\newtheorem{exercise}[subsection]{Exercise}
\newtheorem{situation}[subsection]{Situation}

\theoremstyle{remark}
\newtheorem{remark}[subsection]{Remark}
\newtheorem{remarks}[subsection]{Remarks}

\numberwithin{equation}{subsection}

% Macros
%
\def\lim{\mathop{\mathrm{lim}}\nolimits}
\def\colim{\mathop{\mathrm{colim}}\nolimits}
\def\Spec{\mathop{\mathrm{Spec}}}
\def\Hom{\mathop{\mathrm{Hom}}\nolimits}
\def\Ext{\mathop{\mathrm{Ext}}\nolimits}
\def\SheafHom{\mathop{\mathcal{H}\!\mathit{om}}\nolimits}
\def\SheafExt{\mathop{\mathcal{E}\!\mathit{xt}}\nolimits}
\def\Sch{\mathit{Sch}}
\def\Mor{\mathop{\mathrm{Mor}}\nolimits}
\def\Ob{\mathop{\mathrm{Ob}}\nolimits}
\def\Sh{\mathop{\mathit{Sh}}\nolimits}
\def\NL{\mathop{N\!L}\nolimits}
\def\CH{\mathop{\mathrm{CH}}\nolimits}
\def\proetale{{pro\text{-}\acute{e}tale}}
\def\etale{{\acute{e}tale}}
\def\QCoh{\mathit{QCoh}}
\def\Ker{\mathop{\mathrm{Ker}}}
\def\Im{\mathop{\mathrm{Im}}}
\def\Coker{\mathop{\mathrm{Coker}}}
\def\Coim{\mathop{\mathrm{Coim}}}

% Boxtimes
%
\DeclareMathSymbol{\boxtimes}{\mathbin}{AMSa}{"02}

%
% Macros for moduli stacks/spaces
%
\def\QCohstack{\mathcal{QC}\!\mathit{oh}}
\def\Cohstack{\mathcal{C}\!\mathit{oh}}
\def\Spacesstack{\mathcal{S}\!\mathit{paces}}
\def\Quotfunctor{\mathrm{Quot}}
\def\Hilbfunctor{\mathrm{Hilb}}
\def\Curvesstack{\mathcal{C}\!\mathit{urves}}
\def\Polarizedstack{\mathcal{P}\!\mathit{olarized}}
\def\Complexesstack{\mathcal{C}\!\mathit{omplexes}}
% \Pic is the operator that assigns to X its picard group, usage \Pic(X)
% \Picardstack_{X/B} denotes the Picard stack of X over B
% \Picardfunctor_{X/B} denotes the Picard functor of X over B
\def\Pic{\mathop{\mathrm{Pic}}\nolimits}
\def\Picardstack{\mathcal{P}\!\mathit{ic}}
\def\Picardfunctor{\mathrm{Pic}}
\def\Deformationcategory{\mathcal{D}\!\mathit{ef}}
\usepackage{cleveref}
\crefname{originaldefinition}{original definition}{original definitions}
\Crefname{originaldefinition}{Original definition}{Original definitions}
\makeatletter
\newlabel{categories-definition-essentially-constant-diagram}{{05PU}{source}{}{sourcepointer-1}{}}
\newlabel{categories-definition-essentially-constant-diagram@cref}{{[originaldefinition][0][]05PU}{[1][0][]source}{}{}{}}
\makeatother
\crefname{originallemma}{original lemma}{original lemmas}
\Crefname{originallemma}{Original lemma}{Original lemmas}
\makeatletter
\newlabel{categories-lemma-cofinal-essentially-constant}{{0A1S}{source}{}{sourcepointer-2}{}}
\newlabel{categories-lemma-cofinal-essentially-constant@cref}{{[originallemma][0][]0A1S}{[1][0][]source}{}{}{}}
\makeatother
\title{Essentially constant diagrams in a Karoubian category split off a constant direct summand}
\author{The Stacks Project Authors}
\date{stacks-project:homology.tex:local-snapshot-20261012; GNU FDL 1.2 or later}
\begin{document}
\maketitle
\noindent\textit{Editorial note.} All delivered mathematical text is the author's own, copied from the pinned Stacks Project chapter file homology.tex; no statement or proof step was rewritten, reordered or removed.\par\medskip
\noindent\textit{Editorial note.} The proof is the author's complete argument for the filtered case in both directions; the closing sentence 'The proof of (2) is dual.' is the author's own duality remark on a proof written out in full inside the counted span, preserved verbatim, and no other step of the argument is declined.\par\medskip
\noindent\textit{Editorial note.} Every labelled reference occurring inside the counted spans is carried in assembly.external\_references as a non-load-bearing navigation pointer (same-chapter ones at their definition site, cross-chapter ones at the invocation inside the counted span), so no load-bearing mathematics is replaced by a stub. No bibliography citation occurs inside the counted spans.\par\medskip
\noindent\textit{Editorial note.} The same-name check for this card scanned every staged Stacks chapter of this repository -- all sources/stacks-*/pinned-original-source.tex files, the card's own chapter excluded -- and the labels found are recorded in dedup.same\_name\_cross\_chapter on the card.\par\medskip
\noindent\textit{Editorial note.} The preamble is the chapter's own native preamble excerpt already used for this source: the comment-type reference, slogan and history environments (native preamble lines 8-13), the Xy-pic setup (native preamble lines 14-18) and the native theorem declarations with the mathematical macros (native preamble lines 154-217, the actual last line of the file).\par\medskip
\noindent\textit{Allowed standard prerequisites (editorial).}\begin{itemize}
\item Essentially constant diagrams. Categories, Definition categories-definition-essentially-constant-diagram. Supplies the object X with the maps X -> M\_i and M\_i -> X whose composites are eventually the identity, which the proof uses both to produce the direct summand decomposition and to verify essential constancy.
\item Cofinal subcategories and essential constancy. Categories, Lemma categories-lemma-cofinal-essentially-constant. Allows replacing the index category by the cofinal (respectively initial) subcategory when checking that a diagram is essentially constant.
\item Karoubian categories. Categories, the definition of a Karoubian (pseudo-abelian) category. Supplies the splitting of the idempotent needed to write M\_\{i'\} as a direct sum of the image of X and the kernel of the map to X.
\end{itemize}\medskip
\ifdefined\hypertarget\hypertarget{sourcepointer-1}{}\fi
\noindent\textit{Source reference.} Original definition 05PU: sources/stacks-categories/pinned-original-source.tex, source lines 3145--3145. Navigation pointer invoked at proof line 7713; the counted span contains the pointer only and the cited result is not reproduced inside it, so the pointer is registered as a non-load-bearing navigation stub in the established form for Stacks material. No bibliography citation is involved.\par\medskip
\ifdefined\hypertarget\hypertarget{sourcepointer-2}{}\fi
\noindent\textit{Source reference.} Original lemma 0A1S: sources/stacks-categories/pinned-original-source.tex, source lines 3439--3439. Navigation pointer invoked at proof line 7733; the counted span contains the pointer only and the cited result is not reproduced inside it, so the pointer is registered as a non-load-bearing navigation stub in the established form for Stacks material. No bibliography citation is involved.\par\medskip
\begin{lemma}
\label{lemma-essentially-constant-into-karoubian}
Let $\mathcal{I}$ be a category, let $\mathcal{A}$ be a pre-additive
Karoubian category, and let $M : \mathcal{I} \to \mathcal{A}$ be a diagram.
\begin{enumerate}
\item Assume $\mathcal{I}$ is filtered. The following are equivalent
\begin{enumerate}
\item $M$ is essentially constant,
\item $X = \colim M$ exists and there exists a cofinal filtered subcategory
$\mathcal{I}' \subset \mathcal{I}$ and for $i' \in \Ob(\mathcal{I}')$
a direct sum decomposition $M_{i'} = X_{i'} \oplus Z_{i'}$ such that
$X_{i'}$ maps isomorphically to $X$ and $Z_{i'}$ to zero in $M_{i''}$
for some $i' \to i''$ in $\mathcal{I}'$.
\end{enumerate}
\item Assume $\mathcal{I}$ is cofiltered. The following are equivalent
\begin{enumerate}
\item $M$ is essentially constant,
\item $X = \lim M$ exists and there exists an initial cofiltered subcategory
$\mathcal{I}' \subset \mathcal{I}$ and for $i' \in \Ob(\mathcal{I}')$
a direct sum decomposition $M_{i'} = X_{i'} \oplus Z_{i'}$
such that $X$ maps isomorphically to $X_{i'}$ and $M_{i''} \to Z_{i'}$
is zero for some $i'' \to i'$ in $\mathcal{I}'$.
\end{enumerate}
\end{enumerate}
\end{lemma}

\begin{proof}
Assume (1)(a), i.e., $\mathcal{I}$ is filtered and $M$ is essentially
constant. Let $X = \colim M_i$. Choose $i$ and $X \to M_i$ as in
Categories, Definition
\ref{categories-definition-essentially-constant-diagram}.
Let $\mathcal{I}'$ be the full subcategory consisting of objects
which are the target of a morphism with source $i$.
Suppose $i' \in \Ob(\mathcal{I}')$ and choose a morphism $i \to i'$.
Then $X \to M_i \to M_{i'}$ composed with $M_{i'} \to X$ is the identity
on $X$. As $\mathcal{A}$ is Karoubian, we find a direct summand
decomposition $M_{i'} = X_{i'} \oplus Z_{i'}$, where
$Z_{i'} = \Ker(M_{i'} \to X)$ and $X_{i'}$ maps isomorphically to $X$.
Pick $i \to k$ and $i' \to k$ such that $M_{i'} \to X \to M_i \to M_k$
equals $M_{i'} \to M_k$ as in Categories, Definition
\ref{categories-definition-essentially-constant-diagram}.
Then we see that $M_{i'} \to M_k$ annihilates $Z_{i'}$.
Thus (1)(b) holds.

\medskip\noindent
Assume (1)(b), i.e., $\mathcal{I}$ is filtered and we have
$\mathcal{I}' \subset \mathcal{I}$ and for $i' \in \Ob(\mathcal{I}')$
a direct sum decomposition $M_{i'} = X_{i'} \oplus Z_{i'}$
as stated in the lemma. To see that $M$ is essentially constant
we can replace $\mathcal{I}$ by $\mathcal{I}'$, see
Categories, Lemma \ref{categories-lemma-cofinal-essentially-constant}.
Pick any $i \in \Ob(\mathcal{I})$
and denote $X \to M_i$ the inverse of the isomorphism $X_i \to X$
followed by the inclusion map $X_i \to M_i$. If $j$ is
a second object, then choose $j \to k$ such that $Z_j \to M_k$ is
zero. Since $\mathcal{I}$ is filtered we may also assume there is
a morphism $i \to k$ (after possibly increasing $k$). Then
$M_j \to X \to M_i \to M_k$ and $M_j \to M_k$ both annihilate $Z_j$.
Thus after postcomposing by a morphism $M_k \to M_l$ which annihilates
the summand $Z_k$, we find that $M_j \to X \to M_i \to M_l$ and
$M_j \to M_l$ are equal, i.e., $M$ is essentially constant.

\medskip\noindent
The proof of (2) is dual.
\end{proof}

\end{document}
